\begin{center}\large\bfseries ADDENDUM.\end{center}
After printing, a slight error was noticed in the comparison table at two locations. At Callao it arose from an incorrect longitude in the work cited, and at St. Helena from a calculation error. I take this opportunity to give the results of a revised calculation and also to compare the theory with observations at eight other locations that have since come to my attention. [The corrections have been incorporated in the reprint. For ease of reference, the comparison above already brings together the observations at those eight locations and those at the original 91 locations.]
\SourceLeaf{180}
The Stockholm observations are by \textsc{Rudberg}, the intensity and inclination from 1832 and the declination from 1833 (\textsc{Poggendorff}’s Annalen, volume 37). The observations in Brussels date from 1832. The declination and inclination are by \textsc{Quetelet} (Bulletins de l’Académie de Bruxelles, vol. VI), and the intensity is by \textsc{Rudberg} (the work by \textsc{Sabine} cited above [p. 154]). I owe to \textsc{Sabine}’s kindness in sharing his results the determinations for the other new locations, as well as the determination of the intensity and a more recent observation of the inclination at Callao. The observations at Lerwick and Valentia were made in 1838 by Captain \textsc{James Ross}, and those at Port Etches, Panama and Oahu in 1837 by Captain \textsc{Belcher}. those at Callao in 1838 were made by him as well. Finally, at Montreal the inclination and intensity were observed in 1838 by Major \textsc{Estcourt}. The declination, however, dates from 1834, and the observer is unnamed.

As for the plate illustrating the investigations developed in Article 12, it should be noted here that the skilful lithographer Mr. \textsc{Rittmüller} tried to show the variation in intensity as well. He did so in two ways, by varying both the thickness of the lines and the shading in the spaces between them.

Because printing of the present volume was delayed, it became possible to include two other maps in addition to the map of the values of $V$ [see Article 30]. Readers owe the first of these, which shows the values of the declination \textit{calculated} from the elements or from the tables, to my esteemed friend and co-editor of the \textit{Resultate}. In order to make the intricate arrangement of the system of lines of equal declination as clear as possible, the points where the declination has a maximum value were calculated with particular care, as were those where two lines of equal declination cross (or where one crosses itself). There are two points of the first kind and four of the second. What they have in common is that the first differential of the declination vanishes there in every direction. It hardly needs saying that in regions where the declination changes slowly in all directions, as in southern and southeastern Asia, small changes in the values of the declination can already produce very large changes in the arrangement of the lines.
